% =====================================================================
%  CAPITULO 9 - METODO DE MAXWELL (ANALISE DE MALHAS)
% =====================================================================
\section{Método de Maxwell (Análise de Malhas)}

\begin{conceito}[Antes de começar]
O método de Maxwell (ou análise de malhas) usa a \textbf{lei das malhas} de
Kirchhoff, adotando uma \emph{corrente de malha} que circula em cada laço
fechado. Roteiro:
\begin{enumerate}[leftmargin=1.7em,itemsep=1pt,topsep=2pt]
\item Adote um sentido (geralmente horário) para a corrente de cada malha
      ($I_1, I_2, \dots$).
\item Em cada malha, some as quedas de tensão: um resistor \emph{compartilhado}
      por duas malhas é percorrido pela \emph{diferença} das correntes.
\item Iguale à soma das fem da malha (positivas se a corrente sai do $+$).
\item Resolva o sistema. A corrente real em um resistor compartilhado é a soma
      algébrica das correntes de malha que passam por ele.
\end{enumerate}
\end{conceito}

\legendaniveis

\blocofixacao

\begin{questao}[1]{}
No método das correntes de malha, um resistor compartilhado por duas malhas
adjacentes é percorrido por uma corrente igual a:
\begin{alternativas}
\item soma das duas correntes de malha, sempre.
\item diferença das duas correntes de malha.
\item média das duas correntes de malha.
\item apenas a corrente da primeira malha.
\item zero.
\end{alternativas}
\resposta{(b)}
\resolucao{As duas correntes de malha atravessam o resistor comum em sentidos, em
geral, opostos: a corrente real é a \emph{diferença} algébrica $I_1 - I_2$ (com o
sinal dependendo dos sentidos adotados).}
\end{questao}

\blocoaplicacao

\begin{questao}[2]{}
Aplique o método de Maxwell ao circuito da figura e determine as correntes de
malha $I_1$ e $I_2$.

\circuitoq{%
\begin{circuitikz}[scale=1,transform shape,line width=0.8pt]
 % malha 1 (esquerda)
 \draw (0,0) to[battery1,l_={$\SI{20}{\volt}$}] (0,3)
       to[R,l={$R_1=\SI{4}{\ohm}$}] (3,3) coordinate(top);
 \draw (top) to[R,l={$R_2=\SI{8}{\ohm}$}] (3,0) coordinate(bot) -- (0,0);
 % malha 2 (direita)
 \draw (top) to[R,l={$R_3=\SI{4}{\ohm}$}] (6,3) -- (6,0) -- (bot);
 \draw[->,color=loopcol,line width=1.4pt] (1.4,1.5) arc (0:310:0.5);
 \node[loopcol] at (1.4,1.5) {\footnotesize$I_1$};
 \draw[->,color=loopcol,line width=1.4pt] (4.6,1.5) arc (0:310:0.5);
 \node[loopcol] at (4.6,1.5) {\footnotesize$I_2$};
\end{circuitikz}}

\resposta{$I_1 = \SI{3}{\ampere}$; $I_2 = \SI{2}{\ampere}$}
\resolucao{Correntes de malha no sentido horário. O resistor $R_2$ (comum) é
percorrido por $I_1 - I_2$.\\
\textbf{Malha 1:} $20 = R_1 I_1 + R_2(I_1 - I_2) = 4I_1 + 8(I_1 - I_2)$, ou seja
$12I_1 - 8I_2 = 20$.\\
\textbf{Malha 2:} $0 = R_3 I_2 + R_2(I_2 - I_1) = 4I_2 + 8(I_2 - I_1)$, ou seja
$-8I_1 + 12I_2 = 0 \Rightarrow I_1 = \dfrac{3}{2}I_2$.\\
Substituindo: $12\cdot\dfrac{3}{2}I_2 - 8I_2 = 20 \Rightarrow 10I_2 = 20
\Rightarrow I_2 = \SI{2}{\ampere}$ e $I_1 = \SI{3}{\ampere}$.\\
\textbf{Corrente real em $R_2$:} $I_1 - I_2 = \SI{1}{\ampere}$ (de cima para
baixo).}
\end{questao}

\begin{questao}[2]{}
No circuito de duas malhas da figura, com duas fontes, determine $I_1$ e $I_2$.

\circuitoq{%
\begin{circuitikz}[scale=1,transform shape,line width=0.8pt]
 \draw (0,0) to[battery1,l_={$E_1=\SI{18}{\volt}$}] (0,3)
       to[R,l={$R_1=\SI{2}{\ohm}$}] (3,3) coordinate(top);
 \draw (top) to[R,l={$R_2=\SI{6}{\ohm}$}] (3,0) coordinate(bot) -- (0,0);
 \draw (top) to[R,l={$R_3=\SI{6}{\ohm}$}] (6,3)
       to[battery1,l={$E_2=\SI{6}{\volt}$},invert] (6,0) -- (bot);
 \draw[->,color=loopcol,line width=1.4pt] (1.4,1.5) arc (0:310:0.5);
 \node[loopcol] at (1.4,1.5) {\footnotesize$I_1$};
 \draw[->,color=loopcol,line width=1.4pt] (4.6,1.5) arc (0:310:0.5);
 \node[loopcol] at (4.6,1.5) {\footnotesize$I_2$};
\end{circuitikz}}

\resposta{$I_1 = \SI{3}{\ampere}$; $I_2 = \SI{1}{\ampere}$}
\resolucao{\textbf{Malha 1:} $E_1 = R_1 I_1 + R_2(I_1 - I_2)$:
\[
18 = 2I_1 + 6(I_1 - I_2) \Rightarrow 8I_1 - 6I_2 = 18.
\]
\textbf{Malha 2:} percorrendo com $E_2$ opondo-se a $I_2$:
$-E_2 = R_3 I_2 + R_2(I_2 - I_1)$, isto é
\[
-6 = 6I_2 + 6(I_2 - I_1) \Rightarrow -6I_1 + 12I_2 = -6.
\]
Resolvendo o sistema:
\[
\begin{cases}
8I_1 - 6I_2 = 18\\
-6I_1 + 12I_2 = -6
\end{cases}
\Rightarrow I_1 = \SI{3}{\ampere},\ I_2 = \SI{1}{\ampere}.
\]
A corrente no resistor central $R_2$ é $I_1 - I_2 = \SI{2}{\ampere}$.
\textbf{Verificação} (malha 1): $2\cdot3 + 6\cdot2 = 6+12 = \SI{18}{\volt} = E_1$.
Confere. (Este é o mesmo circuito resolvido por Kirchhoff no Capítulo 2 --- os
dois métodos dão o mesmo resultado.)}
\end{questao}

\begin{questao}[2]{}
No circuito da figura, com uma única fonte, aplique o método de Maxwell e
determine as correntes de malha $I_1$ e $I_2$ e a corrente real no resistor
central.

\circuitoq{%
\begin{circuitikz}[scale=1,transform shape,line width=0.8pt]
 \draw (0,0) to[battery1,l_={$\SI{24}{\volt}$}] (0,3)
       to[R,l={$R_1=\SI{2}{\ohm}$}] (3,3) coordinate(top);
 \draw (top) to[R,l={$R_2=\SI{6}{\ohm}$}] (3,0) coordinate(bot) -- (0,0);
 \draw (top) to[R,l={$R_3=\SI{3}{\ohm}$}] (6,3) -- (6,0) -- (bot);
 \draw[->,color=loopcol,line width=1.4pt] (1.4,1.5) arc (0:310:0.5);
 \node[loopcol] at (1.4,1.5){\footnotesize$I_1$};
 \draw[->,color=loopcol,line width=1.4pt] (4.6,1.5) arc (0:310:0.5);
 \node[loopcol] at (4.6,1.5){\footnotesize$I_2$};
\end{circuitikz}}

\resposta{$I_1 = \SI{6}{\ampere}$, $I_2 = \SI{4}{\ampere}$; $I_{R_2}=\SI{2}{\ampere}$}
\resolucao{Correntes de malha no sentido horário; $R_2$ (comum) conduz
$I_1 - I_2$.\\
\textbf{Malha 1:} $24 = R_1 I_1 + R_2(I_1-I_2) = 2I_1 + 6(I_1-I_2)$, ou
$8I_1 - 6I_2 = 24$.\\
\textbf{Malha 2:} $0 = R_3 I_2 + R_2(I_2-I_1) = 3I_2 + 6(I_2-I_1)$, ou
$-6I_1 + 9I_2 = 0 \Rightarrow I_1 = \tfrac{3}{2}I_2$.\\
Substituindo: $8\cdot\tfrac32 I_2 - 6I_2 = 24 \Rightarrow 6I_2 = 24
\Rightarrow I_2 = \SI{4}{\ampere}$ e $I_1 = \SI{6}{\ampere}$.\\
Corrente real em $R_2$: $I_1 - I_2 = \SI{2}{\ampere}$ (de cima para baixo).}
\end{questao}

\begin{questao}[2]{}
Sintetize, em uma tabela comparativa, quando usar cada técnica de análise
(nodal simples, supernó, malhas simples, supermalha), indicando o "sintoma" que
identifica cada caso.
\resposta{Ver tabela comentada}
\resolucao{
\begin{center}\small\sffamily
\setlength{\arraystretch}{1.35}
\begin{tabular}{@{}p{3.1cm} p{4.6cm} p{4.6cm}@{}}
\toprule
\textbf{Técnica} & \textbf{Sintoma que a indica} & \textbf{Equação extra}\\
\midrule
Nodal simples & Fontes de corrente; fontes de tensão com um terminal no terra
  & --- \\
\textbf{Supernó} & Fonte de tensão \emph{entre dois nós}, nenhum no terra
  & $V_a - V_b = E$ (restrição)\\
Malhas simples & Fontes de tensão; circuito planar
  & --- \\
\textbf{Supermalha} & Fonte de corrente \emph{compartilhada} por duas malhas
  & $I_b - I_a = I_{\text{fonte}}$ (restrição)\\
\bottomrule
\end{tabular}
\end{center}

\medskip
\textbf{O princípio unificador.} Em ambos os "super" métodos, a fonte
problemática é aquela cuja \emph{variável conjugada} é desconhecida: na fonte de
tensão não se sabe a \emph{corrente}; na fonte de corrente não se sabe a
\emph{tensão}. A estratégia é sempre a mesma:
\begin{enumerate}[leftmargin=1.7em,itemsep=1pt,topsep=2pt]
\item escreva a equação de conservação em um contorno que \emph{contorne} o
      elemento problemático (englobando-o, no supernó; evitando-o, na supermalha);
\item recupere a equação que falta a partir da \emph{própria definição} da fonte.
\end{enumerate}
Reconhecer esse padrão evita decorar dois procedimentos distintos --- é um só
princípio aplicado de forma dual.}
\end{questao}

\blocoaprofundamento

\begin{questao}[3]{}
Um estudante deve escolher entre análise nodal e análise de malhas para resolver
um circuito planar com $N$ nós e $M$ malhas. Explique o critério de escolha e
aplique-o ao circuito de duas malhas e três nós da questão anterior.
\resposta{Escolher o método com menor número de equações; aqui, malhas (2 eq.)}
\resolucao{Cada método gera um sistema com tantas equações quantas forem as
incógnitas: a análise \textbf{nodal} produz $(N-1)$ equações (nós menos o terra),
e a análise de \textbf{malhas}, $M$ equações (número de laços independentes).
Escolhe-se o que resulta em \emph{menos} equações.\\
No circuito anterior, há $N = 3$ nós (portanto $N-1 = 2$ equações nodais) e
$M = 2$ malhas (2 equações de malha) --- empate. Como havia \emph{fontes de
tensão}, as malhas costumam ser mais diretas (as fem entram sem manipulação),
então a análise de malhas foi a escolha natural. Se o circuito tivesse muitas
fontes de \emph{corrente} e poucos nós, a análise nodal seria preferível.
\textbf{Regra prática:} conte $N-1$ e $M$; havendo empate, deixe as fontes
decidirem --- fontes de tensão favorecem malhas; fontes de corrente favorecem
nós.}
\end{questao}

\begin{questao}[3]{}
Compare os resultados obtidos para um mesmo circuito de duas malhas resolvido por
três caminhos: (i) leis de Kirchhoff diretas, (ii) método de Maxwell (malhas) e
(iii) análise nodal. O que se pode concluir?
\resposta{Ver comentário}
\resolucao{Os três métodos são \emph{equivalentes} --- todos derivam das mesmas
leis de Kirchhoff (nós e malhas) e da lei de Ohm, e por isso levam
\textbf{exatamente ao mesmo resultado} para correntes e tensões. O que muda é a
\emph{organização} do trabalho:
\begin{itemize}[leftmargin=1.6em,itemsep=1pt,topsep=2pt]
\item \textbf{Kirchhoff direto}: intuitivo, mas gera muitas incógnitas
      (uma corrente por ramo) e equações.
\item \textbf{Maxwell (malhas)}: reduz o número de incógnitas ao número de malhas;
      ideal quando há fontes de \emph{tensão}.
\item \textbf{Nodal}: reduz ao número de nós menos um; ideal quando há fontes de
      \emph{corrente}.
\end{itemize}
A escolha é uma questão de \emph{eficiência}, não de correção. Dominar os três dá
ao estudante a flexibilidade de escolher o mais econômico para cada circuito.}
\end{questao}

\begin{questao}[3]{}
No circuito da figura, uma fonte de corrente de $\SI{2}{\ampere}$ é
\textbf{compartilhada} pelas malhas 1 e 2. Determine as correntes de malha $I_1$ e
$I_2$ pelo método da \textbf{supermalha}.

\circuitoq{%
\begin{circuitikz}[scale=1,transform shape,line width=0.8pt,american]
 % contorno externo
 \draw (0,0) to[battery1,l_={$\SI{24}{\volt}$}] (0,3)
       to[R,l={$R_1=\SI{4}{\ohm}$}] (3,3) coordinate(top);
 % fonte de corrente compartilhada (ramo central)
 \draw (top) to[isource,l={$\SI{2}{\ampere}$}] (3,0) coordinate(bot);
 \draw (bot) -- (0,0);
 % malha 2
 \draw (top) to[R,l={$R_2=\SI{4}{\ohm}$}] (6,3) -- (6,0) -- (bot);
 % setas de malha
 \draw[->,color=loopcol,line width=1.3pt] (1.4,1.5) arc (0:310:0.45);
 \node[loopcol] at (1.4,1.5){\footnotesize$I_1$};
 \draw[->,color=loopcol,line width=1.3pt] (4.7,1.5) arc (0:310:0.45);
 \node[loopcol] at (4.7,1.5){\footnotesize$I_2$};
 % contorno da supermalha
 \draw[dashed,laranja,line width=1pt,rounded corners=10pt]
       (-0.55,-0.5) rectangle (6.55,3.6);
 \node[laranja,font=\footnotesize\sffamily] at (3,-0.85) {supermalha (contorno externo)};
\end{circuitikz}}

\resposta{$I_1 = \SI{2}{\ampere}$; $I_2 = \SI{4}{\ampere}$}
\resolucao{\textbf{Por que o método normal falha:} a tensão sobre a fonte de
corrente é desconhecida, então não se pode escrever a lei das malhas
isoladamente na malha 1 nem na 2.

\medskip
\textbf{Equação 1 --- lei das malhas na supermalha.} Percorremos o
\emph{contorno externo} (tracejado laranja), que \textbf{evita} o ramo da fonte de
corrente. Nesse percurso encontramos a fonte de \SI{24}{\volt}, $R_1$ (percorrido
por $I_1$) e $R_2$ (percorrido por $I_2$):
\[
24 = R_1 I_1 + R_2 I_2 = 4I_1 + 4I_2.
\]

\medskip
\textbf{Equação 2 --- restrição da fonte de corrente.} A fonte no ramo
compartilhado impõe a diferença entre as correntes de malha:
\[
I_2 - I_1 = 2.
\]

\medskip
\textbf{Resolvendo.} Da segunda, $I_2 = I_1+2$. Substituindo:
\[
4I_1 + 4(I_1+2) = 24
\;\Rightarrow\;
8I_1 = 16
\;\Rightarrow\;
I_1 = \SI{2}{\ampere},\quad I_2 = \SI{4}{\ampere}.
\]
\textbf{Verificação:} $4\cdot2 + 4\cdot4 = 8+16 = \SI{24}{\volt}$. Confere.\\
\textbf{Consistência:} a fonte de corrente conduz $I_2 - I_1 = \SI{2}{\ampere}$,
exatamente seu valor nominal.}
\end{questao}

\begin{questao}[3]{}
Para o circuito de \textbf{três malhas} da figura, monte o sistema matricial por
inspeção e resolva-o pela \textbf{regra de Cramer}.

\circuitoq{%
\begin{circuitikz}[scale=0.95,transform shape,line width=0.8pt]
 % malha 1
 \draw (0,0) to[battery1,l_={$E=\SI{24}{\volt}$}] (0,3)
       to[R,l={$R_1=\SI{4}{\ohm}$}] (3,3) coordinate(t1);
 \draw (t1) to[R,l={$R_2=\SI{4}{\ohm}$}] (3,0) coordinate(b1) -- (0,0);
 % malha 2
 \draw (t1) to[R,l={$R_3=\SI{2}{\ohm}$}] (6,3) coordinate(t2);
 \draw (t2) to[R,l={$R_4=\SI{4}{\ohm}$}] (6,0) coordinate(b2) -- (b1);
 % malha 3
 \draw (t2) to[R,l={$R_5=\SI{4}{\ohm}$}] (9,3) -- (9,0) -- (b2);
 % setas
 \draw[->,color=loopcol,line width=1.3pt] (1.5,1.5) arc (0:310:0.42);
 \node[loopcol] at (1.5,1.5){\footnotesize$I_1$};
 \draw[->,color=loopcol,line width=1.3pt] (4.6,1.5) arc (0:310:0.42);
 \node[loopcol] at (4.6,1.5){\footnotesize$I_2$};
 \draw[->,color=loopcol,line width=1.3pt] (7.7,1.5) arc (0:310:0.42);
 \node[loopcol] at (7.7,1.5){\footnotesize$I_3$};
\end{circuitikz}}

\resposta{$I_1=\SI{4}{\ampere}$, $I_2=\SI{2}{\ampere}$, $I_3=\SI{1}{\ampere}$}
\resolucao{\textbf{Passo 1 --- montagem por inspeção.}
\begin{itemize}[leftmargin=1.6em,itemsep=1pt,topsep=2pt]
\item Malha 1: resistências $R_1+R_2 = 4+4 = 8$; compartilha $R_2=4$ com a malha 2;
      não toca a malha 3.
\item Malha 2: $R_2+R_3+R_4 = 4+2+4 = 10$; compartilha $R_2=4$ com a 1 e $R_4=4$
      com a 3.
\item Malha 3: $R_4+R_5 = 4+4 = 8$; compartilha $R_4=4$ com a malha 2.
\end{itemize}
\[
\begin{bmatrix}
8 & -4 & 0\\
-4 & 10 & -4\\
0 & -4 & 8
\end{bmatrix}
\begin{bmatrix}I_1\\I_2\\I_3\end{bmatrix}
=
\begin{bmatrix}24\\0\\0\end{bmatrix}
\]
(Repare na \emph{simetria} da matriz --- primeira verificação de que a montagem
está correta.)

\smallskip
\textbf{Passo 2 --- determinante principal.} Expandindo pela primeira linha:
\[
\Delta = 8\,(10\cdot8 - 16) - (-4)\,(-4\cdot8 - 0)
       = 8(64) - 4(32) = 512 - 128 = 384.
\]

\smallskip
\textbf{Passo 3 --- determinantes das colunas substituídas.}
\[
\Delta_1=\begin{vmatrix}24&-4&0\\0&10&-4\\0&-4&8\end{vmatrix}
=24(80-16)=1536,
\qquad
\Delta_2=\begin{vmatrix}8&24&0\\-4&0&-4\\0&0&8\end{vmatrix}=768,
\]
\[
\Delta_3=\begin{vmatrix}8&-4&24\\-4&10&0\\0&-4&0\end{vmatrix}=384.
\]

\smallskip
\textbf{Passo 4 --- aplicar Cramer.}
\[
I_1=\frac{1536}{384}=\SI{4}{\ampere},\quad
I_2=\frac{768}{384}=\SI{2}{\ampere},\quad
I_3=\frac{384}{384}=\SI{1}{\ampere}.
\]

\smallskip
\textbf{Verificação por balanço de potência} (o teste mais confiável):
\[
P_{\text{fonte}}=E\,I_1=24\cdot4=\SI{96}{\watt},
\]
\[
P_{\text{resistores}} = 4(4)^2 + 4(4-2)^2 + 2(2)^2 + 4(2-1)^2 + 4(1)^2
= 64+16+8+4+4 = \SI{96}{\watt}.
\]
Os valores coincidem: a solução está correta.}
\end{questao}

\begin{questao}[3]{}
No circuito da questão anterior, determine as correntes \emph{reais} nos
resistores compartilhados $R_2$ e $R_4$, e explique por que elas são menores que
as correntes de malha.
\resposta{$I_{R_2}=\SI{2}{\ampere}$; $I_{R_4}=\SI{1}{\ampere}$}
\resolucao{Os resistores compartilhados são percorridos pela \emph{diferença} das
correntes de malha que os atravessam:
\[
I_{R_2}=I_1-I_2 = 4-2 = \SI{2}{\ampere},
\qquad
I_{R_4}=I_2-I_3 = 2-1 = \SI{1}{\ampere}.
\]
Elas são menores porque as duas correntes de malha atravessam o resistor comum em
sentidos \emph{opostos}, cancelando-se parcialmente. As correntes de malha são um
\emph{artifício matemático} --- não são fisicamente mensuráveis; o que um
amperímetro mediria é a corrente de ramo, ou seja, essa diferença.\\
\textbf{Consequência prática:} ao dimensionar um resistor compartilhado, use a
corrente de \emph{ramo} (menor), não a de malha --- caso contrário o
superdimensionamento é desnecessário.}
\end{questao}

\begin{questao}[3]{}
Considere a matriz de malhas de um circuito de quatro malhas:
\[
[\mathbf{R}]=
\begin{bmatrix}
10 & -4 & 0 & -2\\
-4 & 12 & -3 & 0\\
0 & -3 & 9 & -2\\
-2 & 0 & -2 & 8
\end{bmatrix}
\]
\begin{itensq}
\item Sem resolver o sistema, extraia do arranjo a topologia: quais malhas são
      adjacentes e qual o valor de cada resistência compartilhada?
\item Qual a resistência total da malha 3?
\item Por que todos os elementos fora da diagonal são negativos ou nulos?
\end{itensq}
\resposta{Ver leitura estrutural}
\resolucao{(a) Os elementos \emph{fora} da diagonal revelam a topologia. Um valor
não-nulo em $R_{ij}$ indica que as malhas $i$ e $j$ compartilham um resistor de
valor $|R_{ij}|$:
\begin{itemize}[leftmargin=1.6em,itemsep=1pt,topsep=2pt]
\item malhas 1--2: compartilham $\SI{4}{\ohm}$;
\item malhas 1--4: compartilham $\SI{2}{\ohm}$;
\item malhas 2--3: compartilham $\SI{3}{\ohm}$;
\item malhas 3--4: compartilham $\SI{2}{\ohm}$;
\item malhas 1--3 e 2--4: \emph{não são adjacentes} (elementos nulos).
\end{itemize}
A topologia é, portanto, um \textbf{anel} de quatro malhas
($1\!-\!2\!-\!3\!-\!4\!-\!1$).\\
(b) A diagonal dá a soma de todas as resistências da malha:
$R_{33} = \SI{9}{\ohm}$, dos quais \SI{3}{\ohm} são compartilhados com a malha 2 e
\SI{2}{\ohm} com a malha 4; logo, \SI{4}{\ohm} são \emph{exclusivos} da malha 3.\\
(c) Porque as correntes de malha adjacentes percorrem o resistor comum em
sentidos \emph{opostos} (ambas adotadas no mesmo sentido de rotação). A queda de
tensão que a malha vizinha provoca é, assim, subtrativa --- daí o sinal negativo.
\textbf{Isso torna a matriz simétrica e diagonalmente dominante}, o que garante
$\Delta \neq 0$ e a existência de solução única.\\
\textbf{Habilidade desenvolvida:} \emph{ler} um circuito a partir da sua matriz é
tão importante quanto montá-la --- é assim que se depuram erros em simulações e
em softwares de análise (SPICE).}
\end{questao}

\begin{questao}[3]{}
O circuito da figura tem \textbf{três malhas} dispostas em "L" --- a terceira
malha fica \emph{abaixo} da segunda, e não alinhada com as demais. Monte o
sistema matricial e resolva por Cramer.

\circuitoq{%
\begin{circuitikz}[scale=0.95,transform shape,line width=0.8pt]
 % malha 1 (superior esquerda)
 \draw (0,2.6) to[\fonteV,l_={$E=\SI{12}{\volt}$}] (0,5.4)
       to[R,l={$\SI{2}{\ohm}$}] (3,5.4) coordinate(t1);
 \draw (t1) to[R,l={$R_a=\SI{2}{\ohm}$}] (3,2.6) coordinate(m1) -- (0,2.6);
 % malha 2 (superior direita)
 \draw (t1) to[R,l={$\SI{1}{\ohm}$}] (6,5.4) coordinate(t2);
 \draw (t2) -- (6,2.6) coordinate(m2) -- (m1);
 % malha 3 (INFERIOR, abaixo da malha 1)
 \draw (m1) to[R,l_={$R_b=\SI{2}{\ohm}$}] (3,0) coordinate(b1);
 \draw (b1) to[R,l_={$\SI{4}{\ohm}$}] (6,0) coordinate(b2) -- (m2);
 % setas de malha
 \draw[->,color=loopcol,line width=1.3pt] (1.5,4.0) arc (0:310:0.42);
 \node[loopcol] at (1.5,4.0){\footnotesize$I_1$};
 \draw[->,color=loopcol,line width=1.3pt] (4.5,4.0) arc (0:310:0.42);
 \node[loopcol] at (4.5,4.0){\footnotesize$I_2$};
 \draw[->,color=loopcol,line width=1.3pt] (4.5,1.3) arc (0:310:0.42);
 \node[loopcol] at (4.5,1.3){\footnotesize$I_3$};
\end{circuitikz}}

\resposta{$I_1=\SI{3}{\ampere}$, $I_2=\SI{2}{\ampere}$, $I_3=\SI{1}{\ampere}$}
\resolucao{\textbf{Passo 1 --- identificar os acoplamentos.} Apesar da disposição
em "L", a regra é a mesma: o que importa é \emph{quais malhas compartilham
resistores}.
\begin{itemize}[leftmargin=1.6em,itemsep=1pt,topsep=2pt]
\item Malhas 1 e 2 compartilham $R_a = \SI{2}{\ohm}$;
\item Malhas 1 e 3 compartilham $R_b = \SI{2}{\ohm}$;
\item Malhas 2 e 3 \textbf{não} são adjacentes (elemento nulo na matriz).
\end{itemize}
\textbf{Atenção:} a posição \emph{geométrica} das malhas no desenho é irrelevante
--- o que define a matriz é a \emph{topologia} (quem toca quem).

\textbf{Passo 2 --- montar a matriz.}
\[
\begin{bmatrix}
 6 & -2 & -2\\
-2 &  3 &  0\\
-2 &  0 &  6
\end{bmatrix}
\begin{bmatrix}I_1\\I_2\\I_3\end{bmatrix}
=
\begin{bmatrix}12\\0\\0\end{bmatrix}
\]
Diagonal: $R_{11}=2+2+2=6$; $R_{22}=2+1=3$; $R_{33}=2+4=6$.
Os zeros em $R_{23}=R_{32}=0$ confirmam a não-adjacência entre as malhas 2 e 3.

\textbf{Passo 3 --- determinantes e Cramer.}
\[
\Delta = 6(3\cdot6-0)-(-2)\big[(-2)(6)-0\big]+(-2)\big[0-(-2)(3)\big]
= 108-24-12 = 72.
\]
\[
\Delta_1 = 216,\qquad \Delta_2 = 144,\qquad \Delta_3 = 72,
\]
\[
I_1=\frac{216}{72}=\SI{3}{\ampere},\quad
I_2=\frac{144}{72}=\SI{2}{\ampere},\quad
I_3=\frac{72}{72}=\SI{1}{\ampere}.
\]
\textbf{Verificação rápida pelas linhas 2 e 3:}
$3I_2 = 2I_1 \Rightarrow I_2=\tfrac23\cdot3=2$ \checkmark;
$6I_3 = 2I_1 \Rightarrow I_3=\tfrac13\cdot3=1$ \checkmark.

\textbf{Passo 4 --- balanço de potência.}
\[
P_{\text{res}} = 2(3)^2+2(3-2)^2+2(3-1)^2+1(2)^2+4(1)^2
= 18+2+8+4+4 = \SI{36}{\watt},
\]
\[
P_{\text{fonte}} = 12\cdot3 = \SI{36}{\watt}.\ \checkmark
\]
\textbf{Lição estrutural.} Uma matriz de malhas com \emph{zeros fora da diagonal}
indica malhas não-adjacentes. Quanto mais zeros, mais "esparsa" a matriz e mais
fácil a solução --- e é justamente por isso que programas de simulação (SPICE)
usam algoritmos para matrizes esparsas: circuitos reais têm pouquíssimos
acoplamentos por malha.}
\end{questao}

\blocodesafio

\begin{questao}[4]{}
Um circuito planar tem $b$ ramos, $n$ nós e $m$ malhas independentes.
\begin{itensq}
\item Demonstre a relação $m = b - n + 1$ (teorema de Euler para grafos planares).
\item Um circuito tem 12 ramos e 7 nós. Quantas equações de malha e quantas
      equações nodais seriam necessárias? Qual método é mais econômico?
\end{itensq}
\resposta{(a) demonstração; (b) 6 de malha vs 6 nodais --- empate}
\resolucao{(a) Cada ramo carrega uma corrente desconhecida: são $b$ incógnitas.
A lei dos nós fornece $(n-1)$ equações independentes (a $n$-ésima é combinação
linear das demais). Restam, portanto,
\[
m = b - (n-1) = b - n + 1
\]
equações independentes a serem obtidas pela lei das malhas. Essa é exatamente a
relação de Euler aplicada a grafos planares conexos.\\
(b) Com $b=12$ e $n=7$:
\[
m = 12 - 7 + 1 = 6 \text{ equações de malha};
\qquad
n-1 = 6 \text{ equações nodais}.
\]
Há \textbf{empate} (6 equações em ambos). O critério de desempate passa a ser o
tipo de fonte: predominando fontes de \emph{tensão}, prefira malhas; predominando
fontes de \emph{corrente}, prefira nodal. Se houver muitas fontes de tensão
aterradas, a nodal ganha, pois cada uma delas \emph{elimina} uma incógnita
(o potencial daquele nó é conhecido).}
\end{questao}

\begin{questao}[4]{}
O circuito da figura possui \textbf{três malhas} e \textbf{duas fontes}. Monte o
sistema matricial, resolva-o pela regra de Cramer e interprete o sinal obtido
para $I_3$. Verifique o resultado pelo balanço de potência.

\circuitoq{%
\begin{circuitikz}[scale=0.95,transform shape,line width=0.8pt]
 \draw (0,0) to[\fonteV,l={$E_1=\SI{24}{\volt}$}] (0,3)
       to[R,l={$\SI{2}{\ohm}$}] (3,3) coordinate(t1);
 \draw (t1) to[R,l={$\SI{2}{\ohm}$}] (3,0) coordinate(b1) -- (0,0);
 \draw (t1) to[R,l={$\SI{2}{\ohm}$}] (6,3) coordinate(t2);
 \draw (t2) to[R,l={$\SI{2}{\ohm}$}] (6,0) coordinate(b2) -- (b1);
 \draw (t2) to[R,l={$\SI{6}{\ohm}$}] (9,3)
       to[\fonteV,l={$E_2=\SI{12}{\volt}$},invert] (9,0) -- (b2);
 \foreach \x/\n in {1.5/1,4.5/2,7.5/3}{
   \draw[->,color=loopcol,line width=1.3pt] (\x,1.5) arc (0:310:0.42);
   \node[loopcol] at (\x,1.5){\footnotesize$I_{\n}$};}
\end{circuitikz}}

\resposta{$I_1=\SI{7}{\ampere}$, $I_2=\SI{2}{\ampere}$, $I_3=\SI{-1}{\ampere}$
(sentido oposto ao adotado)}
\resolucao{\textbf{Passo 1 --- identificar os acoplamentos.} Apesar da disposição
em "L", a regra é a mesma: o que importa é \emph{quais malhas compartilham
resistores}.
\begin{itemize}[leftmargin=1.6em,itemsep=1pt,topsep=2pt]
\item Malhas 1 e 2 compartilham $R_a = \SI{4}{\ohm}$;
\item Malhas 1 e 3 compartilham $R_b = \SI{4}{\ohm}$;
\item Malhas 2 e 3 \textbf{não} são adjacentes (elemento nulo na matriz).
\end{itemize}
\textbf{Atenção:} a posição \emph{geométrica} das malhas no desenho é irrelevante
--- o que define a matriz é a \emph{topologia} (quem toca quem).

\textbf{Passo 2 --- montar a matriz.}
\[
\begin{bmatrix}
10 & -4 & -4\\
-4 & 6 & 0\\
-4 & 0 & 8
\end{bmatrix}
\begin{bmatrix}I_1\\I_2\\I_3\end{bmatrix}
=
\begin{bmatrix}24\\0\\0\end{bmatrix}
\]
Diagonal: $R_{11}=2+4+4=10$; $R_{22}=4+2=6$; $R_{33}=4+2+2=8$.
O zero em $R_{23}$ confirma a não-adjacência.

\textbf{Passo 3 --- Cramer.}
\[
\Delta = 10(48-0)-(-4)(-32-0)+(-4)(0+24)=480-128-96=\SI{256}{}.
\]
\[
I_1=\frac{\Delta_1}{\Delta},\quad
I_2=\frac{\Delta_2}{\Delta},\quad
I_3=\frac{\Delta_3}{\Delta}.
\]
Resolvendo o sistema, obtém-se
\[
I_1 = \SI{4}{\ampere},\qquad
I_2 = \SI{2.67}{\ampere},\qquad
I_3 = \SI{2}{\ampere}.
\]
\textbf{Verificação pela malha 2:} $6I_2 = 4I_1 \Rightarrow I_2 = \tfrac{2}{3}I_1$
\checkmark; \textbf{malha 3:} $8I_3 = 4I_1 \Rightarrow I_3 = \tfrac12 I_1$
\checkmark.

\textbf{Lição estrutural.} Uma matriz de malhas com \emph{zeros fora da diagonal}
indica malhas não-adjacentes. Quanto mais zeros, mais "esparsa" a matriz e mais
fácil a solução --- e é justamente por isso que programas de simulação (SPICE)
usam algoritmos para matrizes esparsas: circuitos reais têm pouquíssimos
acoplamentos por malha.}
\end{questao}

\begin{questao}[4]{}
O circuito da figura possui \textbf{quatro malhas} em arranjo $2\times2$.
\begin{itensq}
\item Monte a matriz de malhas por inspeção e verifique sua simetria.
\item Resolva o sistema pela regra de Cramer.
\item Determine as correntes reais em todos os resistores compartilhados.
\item Valide a solução pelo balanço de potência.
\end{itensq}

\circuitoq{%
\begin{circuitikz}[scale=0.95,transform shape,line width=0.8pt]
 % nos da grade 2x2
 % linha superior y=5.2 ; meio y=2.6 ; inferior y=0
 % colunas x=0, 3.4, 6.8
 \draw (0,2.6) to[\fonteV,l_={$E=\SI{48}{\volt}$}] (0,5.2)
       to[R,l={$\SI{4}{\ohm}$}] (3.4,5.2) coordinate(A);
 \draw (A) to[R,l={$\SI{2}{\ohm}$}] (6.8,5.2) coordinate(B);
 % ramo vertical central superior (compartilhado 1-2)
 \draw (A) to[R,l={$\SI{2}{\ohm}$}] (3.4,2.6) coordinate(C);
 % ramo vertical direito (malha 2)
 \draw (B) to[R,l={$\SI{2}{\ohm}$}] (6.8,2.6) coordinate(D);
 % linha do meio
 \draw (0,2.6) -- (C); \draw (C) -- (D);
 % ramo vertical esquerdo inferior (compartilhado 1-3)
 \draw (0,2.6) to[R,l_={$\SI{1}{\ohm}$}] (0,0) coordinate(E);
 % linha inferior
 \draw (E) to[R,l_={$\SI{2}{\ohm}$}] (3.4,0) coordinate(F);
 \draw (F) to[R,l_={$\SI{6}{\ohm}$}] (6.8,0) coordinate(G);
 % ramo vertical central inferior (compartilhado 3-4)
 \draw (C) to[R,l={$\SI{2}{\ohm}$}] (F);
 \draw (D) -- (G);
 % setas
 \draw[->,color=loopcol,line width=1.3pt] (1.7,3.9) arc (0:310:0.4);
 \node[loopcol] at (1.7,3.9){\footnotesize$I_1$};
 \draw[->,color=loopcol,line width=1.3pt] (5.1,3.9) arc (0:310:0.4);
 \node[loopcol] at (5.1,3.9){\footnotesize$I_2$};
 \draw[->,color=loopcol,line width=1.3pt] (1.7,1.3) arc (0:310:0.4);
 \node[loopcol] at (1.7,1.3){\footnotesize$I_3$};
 \draw[->,color=loopcol,line width=1.3pt] (5.1,1.3) arc (0:310:0.4);
 \node[loopcol] at (5.1,1.3){\footnotesize$I_4$};
\end{circuitikz}}

\resposta{$I_1=\SI{8}{\ampere}$, $I_2=\SI{3}{\ampere}$, $I_3=\SI{2}{\ampere}$,
$I_4=\SI{1}{\ampere}$}
\resolucao{\textbf{(a) Montagem por inspeção.} Adjacências: $1$--$2$
($\SI{2}{\ohm}$), $1$--$3$ ($\SI{1}{\ohm}$), $2$--$4$ ($\SI{2}{\ohm}$),
$3$--$4$ ($\SI{2}{\ohm}$); as malhas $1$--$4$ e $2$--$3$ são \emph{diagonais},
logo não-adjacentes (zeros).
\[
[\mathbf R]=
\begin{bmatrix}
 7 & -2 & -1 & 0\\
-2 &  6 &  0 & -2\\
-1 &  0 &  5 & -2\\
 0 & -2 & -2 & 10
\end{bmatrix},
\qquad
[\mathbf E]=
\begin{bmatrix}48\\0\\0\\0\end{bmatrix}
\]
A matriz é \textbf{simétrica} \checkmark\ --- primeira verificação de que a
montagem está correta.

\textbf{(b) Regra de Cramer.}
\[
\Delta=\det[\mathbf R]=1536,
\]
\[
\Delta_1=12288,\quad \Delta_2=4608,\quad \Delta_3=3072,\quad \Delta_4=1536.
\]
Portanto
\[
I_1=\frac{12288}{1536}=\SI{8}{\ampere},\quad
I_2=\frac{4608}{1536}=\SI{3}{\ampere},
\]
\[
I_3=\frac{3072}{1536}=\SI{2}{\ampere},\quad
I_4=\frac{1536}{1536}=\SI{1}{\ampere}.
\]

\textbf{(c) Correntes de ramo} (diferença entre as malhas adjacentes):
\begin{center}\small\sffamily
\begin{tabular}{@{}l c c@{}}
\toprule
\textbf{Resistor} & \textbf{Malhas} & \textbf{Corrente real}\\
\midrule
\SI{2}{\ohm} (central sup.) & 1--2 & $I_1-I_2=\SI{5}{\ampere}$\\
\SI{1}{\ohm} (esquerdo)     & 1--3 & $I_1-I_3=\SI{6}{\ampere}$\\
\SI{2}{\ohm} (direito)      & 2--4 & $I_2-I_4=\SI{2}{\ampere}$\\
\SI{2}{\ohm} (central inf.) & 3--4 & $I_3-I_4=\SI{1}{\ampere}$\\
\bottomrule
\end{tabular}
\end{center}

\textbf{(d) Balanço de potência.}
\[
P_{\text{res}} = 4(8)^2+2(5)^2+1(6)^2+2(3)^2+2(2)^2+2(2)^2+2(1)^2+6(1)^2
\]
\[
= 256+50+36+18+8+8+2+6 = \SI{384}{\watt}.
\]
\[
P_{\text{fonte}} = E\,I_1 = 48\cdot8 = \SI{384}{\watt}.\ \checkmark
\]
\textbf{Comentário sobre o método.} Com $n=4$ malhas, resolver "na mão" por
substituição seria trabalhoso e propenso a erro. A montagem matricial por
inspeção elimina a etapa de escrever equações uma a uma, e o balanço de potência
oferece uma verificação \emph{independente} do resultado --- se ele fecha, a
probabilidade de erro é mínima. \emph{Este é o fluxo de trabalho profissional:
montar sistematicamente, resolver com álgebra linear, validar por energia.}}
\end{questao}

% BEGIN BANCO COMPLEMENTAR Método de Maxwell (Análise de Malhas)
% =====================================================================
% BANCO COMPLEMENTAR --- ANALISE DE MALHAS / METODO DE MAXWELL
% Revisao 2026-08-15
% 7 N1 + 8 N2 + 5 N3 + 5 N4 = 25 questoes.
% N2/N3/N4: todas as questoes partem de CIRCUITOS diferentes e as setas
% de corrente de malha sao desenhadas explicitamente no sentido HORARIO.
% Fontes dependentes permanecem nos bancos especificos, evitando transformar
% toda a secao em repeticoes do mesmo tema.
% =====================================================================

\subsubsection*{Banco complementar --- Método de Maxwell (análise de malhas)}

\begin{atencao}[Convenção gráfica]
Em todos os circuitos deste banco, $I_1,I_2,\ldots$ são adotadas no
\textbf{sentido horário}. Uma corrente calculada negativa significa apenas que
a corrente real daquela malha circula no sentido oposto ao indicado. Os níveis
2, 3 e 4 usam topologias e objetivos diferentes: malhas triangulares, fontes em
ramos compartilhados, fontes de corrente periféricas, redes de três malhas,
ponte resistiva, síntese, máxima potência, identificação de parâmetros e rede
em roda de quatro malhas.
\end{atencao}

\providecommand{\fvMeshCW}[3]{%
  \draw[->,loopcol,line width=1.1pt] (#1)
    arc[start angle=0,end angle=-305,radius=#2cm];%
  \node[loopcol,fill=white,inner sep=1pt] at (#1){\footnotesize #3};}

% =====================================================================
% N1
% =====================================================================
\blocofixacao

\begin{questao}[1]{}
No desenho abaixo, a corrente de malha segue a convenção usada neste capítulo.
Identifique o sentido de $I_1$.
\circuitoq{%
\begin{circuitikz}[american,scale=.9,transform shape]
\draw (0,0) to[\fonteV,l^={$E$}] (0,2.8) to[R,l={$R$}] (4,2.8) -- (4,0)--(0,0);
\fvMeshCW{2.0,1.25}{.42}{$I_1$}
\end{circuitikz}}
\resposta{Horário}
\resolucao{A seta percorre o laço no sentido horário. Essa é apenas uma
convenção de referência; se o resultado algébrico de $I_1$ for negativo, a
corrente real circula no sentido anti-horário.}
\end{questao}

\begin{questao}[1]{}
Duas malhas horárias compartilham um resistor $R_c$. Qual é a corrente real no
resistor, tomando como positivo o sentido de $I_1$ no ramo comum?
\resposta{$i_c=I_1-I_2$}
\resolucao{As correntes horárias atravessam o ramo compartilhado em sentidos
opostos; por isso a corrente de ramo é a diferença algébrica.}
\end{questao}

\begin{questao}[1]{}
Um circuito planar conectado possui $b=11$ ramos e $n=7$ nós. Quantas correntes
de malha independentes são necessárias?
\resposta{$m=5$}
\resolucao{$m=b-n+1=11-7+1=5$.}
\end{questao}

\begin{questao}[1]{}
Uma fonte ideal de corrente está na periferia de uma única malha e sua seta
coincide com a corrente horária $I_2$. O que pode ser escrito imediatamente?
\resposta{$I_2=I_s$}
\resolucao{A fonte pertence somente àquela malha; portanto ela fixa diretamente
a corrente de malha e elimina uma incógnita do sistema.}
\end{questao}

\begin{questao}[1]{}
Uma fonte ideal de corrente está no ramo comum a duas malhas. Qual técnica deve
ser usada?
\resposta{Supermalha, mais a equação de restrição da fonte}
\resolucao{A tensão sobre uma fonte ideal de corrente é desconhecida; por isso
não se escreve LKT separadamente nas duas malhas. Contorna-se o ramo da fonte e
acrescenta-se a relação entre $I_1$ e $I_2$.}
\end{questao}

\begin{questao}[1]{}
Na matriz de uma rede resistiva com correntes de malha todas horárias, o que
representam $R_{kk}$ e $R_{jk}$, $j\ne k$?
\resposta{$R_{kk}$ é a soma das resistências da malha $k$; $R_{jk}$ é o negativo
da resistência compartilhada}
\resolucao{A diagonal reúne todos os resistores percorridos por $I_k$. Fora da
diagonal aparece $-R_c$ porque as correntes horárias percorrem o ramo comum em
sentidos opostos.}
\end{questao}

\begin{questao}[1]{}
Ao resolver uma malha adotada no sentido horário, obtém-se
$I_3=-\SI{0.8}{\ampere}$. Interprete fisicamente.
\resposta{A corrente real tem módulo $\SI{0.8}{\ampere}$ e circula no sentido anti-horário}
\resolucao{O sinal negativo não indica erro: ele informa que a escolha inicial
do sentido foi oposta à corrente física.}
\end{questao}

% =====================================================================
% N2 --- circuitos distintos
% =====================================================================
\blocoaplicacao

\begin{questao}[2]{}
Duas malhas triangulares compartilham o resistor diagonal de
$\SI{6}{\ohm}$. Determine $I_1$, $I_2$ e a corrente no ramo compartilhado.
\circuitoq{%
\begin{circuitikz}[american,scale=.88,transform shape,line width=.8pt]
\coordinate (A) at (0,0); \coordinate (B) at (3.4,4.2); \coordinate (C) at (6.8,0);
\draw (A) to[\fonteV,l^={$\SI{18}{\volt}$}] (0,2.0)
      to[R,l={$\SI{4}{\ohm}$}] (B);
\draw (B) to[R,l={$\SI{3}{\ohm}$}] (5.2,2.1)
      to[R,l={$\SI{3}{\ohm}$}] (C);
\draw (C)--(A);
\draw (B) to[R,l={$\SI{6}{\ohm}$}] (3.4,0)--(A);
\fvMeshCW{1.75,1.45}{.38}{$I_1$}
\fvMeshCW{5.05,1.45}{.38}{$I_2$}
\end{circuitikz}}
\resposta{$I_1=\dfrac{18}{7}\,\si{\ampere}$; $I_2=\dfrac{9}{7}\,\si{\ampere}$;
$i_c=\dfrac{9}{7}\,\si{\ampere}$}
\resolucao{As equações são
\[
10I_1-6I_2=18,\qquad -6I_1+12I_2=0.
\]
Da segunda, $I_1=2I_2$. Logo $14I_2=18$, resultando
$I_2=9/7$ A e $I_1=18/7$ A. No ramo comum,
$i_c=I_1-I_2=9/7$ A.}
\end{questao}

\begin{questao}[2]{}
A fonte de corrente de $\SI{1.5}{\ampere}$ está na periferia da malha 2 e sua
seta coincide com $I_2$. Determine $I_1$, $I_2$ e a corrente no resistor
central de $\SI{5}{\ohm}$.
\circuitoq{%
\begin{circuitikz}[american,scale=.92,transform shape,line width=.8pt]
\draw (0,0) to[\fonteV,l^={$\SI{30}{\volt}$}] (0,3.2)
      to[R,l={$\SI{5}{\ohm}$}] (3.2,3.2) coordinate(t);
\draw (t) to[R,l_={$\SI{5}{\ohm}$}] (3.2,0) coordinate(b)--(0,0);
\draw (t) to[R,l={$\SI{10}{\ohm}$}] (6.4,3.2)
      to[isource,l={$\SI{1.5}{\ampere}$}] (6.4,0)--(b);
\fvMeshCW{1.55,1.05}{.40}{$I_1$}
\fvMeshCW{4.85,1.05}{.40}{$I_2$}
\end{circuitikz}}
\resposta{$I_1=\SI{3.75}{\ampere}$; $I_2=\SI{1.5}{\ampere}$;
$i_c=\SI{2.25}{\ampere}$}
\resolucao{A fonte periférica fixa $I_2=1{,}5$ A. Para a malha 1,
\[
5I_1+5(I_1-I_2)=30,
\]
portanto $10I_1-7{,}5=30$ e $I_1=3{,}75$ A. Assim
$i_c=I_1-I_2=2{,}25$ A.}
\end{questao}

\begin{questao}[2]{}
O ramo comum contém um resistor de $\SI{4}{\ohm}$ em série com uma fonte de
$\SI{6}{\volt}$, com terminal positivo no topo. Determine $I_1$ e $I_2$.
\circuitoq{%
\begin{circuitikz}[american,scale=.9,transform shape,line width=.8pt]
\draw (0,0) to[\fonteV,l^={$\SI{24}{\volt}$}] (0,3.6)
      to[R,l={$\SI{4}{\ohm}$}] (3.4,3.6) coordinate(t);
\draw (t) to[R,l_={$\SI{4}{\ohm}$}] (3.4,2.0)
      to[\fonteV,l_={$\SI{6}{\volt}$}] (3.4,0) coordinate(b)--(0,0);
\draw (t) to[R,l={$\SI{8}{\ohm}$}] (6.8,3.6)--(6.8,0)--(b);
\fvMeshCW{1.65,1.15}{.40}{$I_1$}
\fvMeshCW{5.15,1.15}{.40}{$I_2$}
\end{circuitikz}}
\resposta{$I_1=\SI{3}{\ampere}$; $I_2=\SI{1.5}{\ampere}$}
\resolucao{Percorrendo as duas malhas no sentido horário:
\[
8I_1-4I_2=18,
\qquad
-4I_1+12I_2=6.
\]
A solução é $I_1=3$ A e $I_2=1{,}5$ A. A fonte compartilhada entra com sinais
opostos nas duas equações porque é atravessada em sentidos opostos.}
\end{questao}

\begin{questao}[2]{}
A rede possui três malhas triangulares em torno do nó central $O$. Monte a
matriz por inspeção e determine as três correntes.
\circuitoq{%
\begin{circuitikz}[american,scale=.78,transform shape,line width=.8pt]
\coordinate (A) at (0,0); \coordinate (B) at (4,6); \coordinate (C) at (8,0); \coordinate (O) at (4,2.1);
% borda AB: R=4 e E1=16
\draw (A) to[\fonteV,l^={$\SI{16}{\volt}$}] (1.15,1.75)
      to[R,l={$\SI{4}{\ohm}$}] (B);
% borda BC: R=5 e E2=11
\draw (B) to[R,l={$\SI{5}{\ohm}$}] (6.85,1.75)
      to[\fonteV,l={$\SI{11}{\volt}$}] (C);
% borda CA: R=6 e E3=1
\draw (C) to[R,l_={$\SI{6}{\ohm}$}] (4.7,0)
      to[\fonteV,l_={$\SI{1}{\volt}$}] (A);
% raios compartilhados
\draw (B) to[R,l={$\SI{2}{\ohm}$}] (O);
\draw (C) to[R,l={$\SI{3}{\ohm}$}] (O);
\draw (A) to[R,l_={$\SI{1}{\ohm}$}] (O);
\fvMeshCW{2.8,3.5}{.34}{$I_1$}
\fvMeshCW{5.4,3.5}{.34}{$I_2$}
\fvMeshCW{4.0,.95}{.34}{$I_3$}
\node[fill=white,inner sep=1pt] at(O){$O$};
\end{circuitikz}}
\resposta{$I_1=\SI{3}{\ampere}$; $I_2=\SI{2}{\ampere}$; $I_3=\SI{1}{\ampere}$}
\resolucao{Por inspeção,
\[
\begin{bmatrix}
7&-2&-1\\
-2&10&-3\\
-1&-3&10
\end{bmatrix}
\begin{bmatrix}I_1\\I_2\\I_3\end{bmatrix}
=
\begin{bmatrix}16\\11\\1\end{bmatrix}.
\]
Substituindo $(3,2,1)$ verifica-se cada linha; logo essa é a solução. O desenho
mostra que a matriz depende de quais raios são compartilhados, e não de uma
sequência linear de malhas.}
\end{questao}

\begin{questao}[2]{}
As duas fontes do circuito atuam em oposição. Determine as correntes horárias e
interprete o sinal de $I_2$.
\circuitoq{%
\begin{circuitikz}[american,scale=.92,transform shape,line width=.8pt]
\draw (0,0) to[\fonteV,l^={$\SI{20}{\volt}$}] (0,3.2)
      to[R,l={$\SI{6}{\ohm}$}] (3.2,3.2) coordinate(t);
\draw (t) to[R,l_={$\SI{4}{\ohm}$}] (3.2,0) coordinate(b)--(0,0);
\draw (t) to[R,l={$\SI{6}{\ohm}$}] (6.4,3.2)
      to[\fonteV,l_={$\SI{12}{\volt}$}] (6.4,0)--(b);
\fvMeshCW{1.55,1.05}{.40}{$I_1$}
\fvMeshCW{4.85,1.05}{.40}{$I_2$}
\end{circuitikz}}
\resposta{$I_1=\dfrac{38}{21}\,\si{\ampere}$;
$I_2=-\dfrac{10}{21}\,\si{\ampere}$}
\resolucao{O sistema é
\[
\begin{bmatrix}10&-4\\-4&10\end{bmatrix}
\begin{bmatrix}I_1\\I_2\end{bmatrix}
=
\begin{bmatrix}20\\-12\end{bmatrix}.
\]
Daí $I_1=38/21$ A e $I_2=-10/21$ A. Portanto a corrente física da segunda
malha tem módulo $10/21$ A e circula no sentido anti-horário.}
\end{questao}

\begin{questao}[2]{}
Escolha o valor e a polaridade da fonte $E$ para que a corrente no resistor
compartilhado de $\SI{4}{\ohm}$ seja nula.
\circuitoq{%
\begin{circuitikz}[american,scale=.92,transform shape,line width=.8pt]
\draw (0,0) to[\fonteV,l^={$\SI{24}{\volt}$}] (0,3.2)
      to[R,l={$\SI{4}{\ohm}$}] (3.2,3.2) coordinate(t);
\draw (t) to[R,l_={$\SI{4}{\ohm}$}] (3.2,0) coordinate(b)--(0,0);
\draw (t) to[R,l={$\SI{8}{\ohm}$}] (6.4,3.2)
      to[\fonteV,l_={$E$}] (6.4,0)--(b);
\fvMeshCW{1.55,1.05}{.40}{$I_1$}
\fvMeshCW{4.85,1.05}{.40}{$I_2$}
\end{circuitikz}}
\resposta{$E=\SI{48}{\volt}$, impulsionando $I_2$ no sentido horário}
\resolucao{Corrente nula no ramo comum exige $I_1=I_2=I$. A primeira malha dá
$(8-4)I=24$, ou $4I=24$, logo $I=6$ A. Na segunda,
$(12-4)I=E$, portanto $E=8\cdot6=48$ V, com polaridade que impulsione $I_2$.}
\end{questao}

\begin{questao}[2]{}
Determine as correntes de malha e faça o balanço de potência.
\circuitoq{%
\begin{circuitikz}[american,scale=.92,transform shape,line width=.8pt]
\draw (0,0) to[\fonteV,l^={$\SI{36}{\volt}$}] (0,3.2)
      to[R,l={$\SI{6}{\ohm}$}] (3.2,3.2) coordinate(t);
\draw (t) to[R,l_={$\SI{3}{\ohm}$}] (3.2,0) coordinate(b)--(0,0);
\draw (t) to[R,l={$\SI{6}{\ohm}$}] (6.4,3.2)
      to[\fonteV,l_={$\SI{12}{\volt}$},invert] (6.4,0)--(b);
\fvMeshCW{1.55,1.05}{.40}{$I_1$}
\fvMeshCW{4.85,1.05}{.40}{$I_2$}
\end{circuitikz}}
\resposta{$I_1=\SI{5}{\ampere}$; $I_2=\SI{3}{\ampere}$;
$P_{fontes}=P_R=\SI{216}{\watt}$}
\resolucao{O sistema
\[
\begin{bmatrix}9&-3\\-3&9\end{bmatrix}
\begin{bmatrix}I_1\\I_2\end{bmatrix}
=
\begin{bmatrix}36\\12\end{bmatrix}
\]
fornece $(I_1,I_2)=(5,3)$ A. A corrente no resistor comum é $2$ A. Assim
\[
P_R=6(5)^2+6(3)^2+3(2)^2=150+54+12=216\,\text{W}.
\]
As fontes fornecem $36(5)+12(3)=216$ W.}
\end{questao}

\begin{questao}[2]{}
A fonte de corrente da malha 2 fixa $I_2=\SI{2}{\ampere}$. Determine $I_1$ e a
tensão $v_s$ da fonte de corrente, definida positiva do terminal superior para
o inferior.
\circuitoq{%
\begin{circuitikz}[american,scale=.92,transform shape,line width=.8pt]
\draw (0,0) to[\fonteV,l^={$\SI{20}{\volt}$}] (0,3.2)
      to[R,l={$\SI{5}{\ohm}$}] (3.2,3.2) coordinate(t);
\draw (t) to[R,l_={$\SI{5}{\ohm}$}] (3.2,0) coordinate(b)--(0,0);
\draw (t) to[R,l={$\SI{10}{\ohm}$}] (6.4,3.2)
      to[isource,l={$\SI{2}{\ampere}$}] (6.4,0)--(b);
\node[right] at(6.75,1.6){$v_s$};
\fvMeshCW{1.55,1.05}{.40}{$I_1$}
\fvMeshCW{4.85,1.05}{.40}{$I_2$}
\end{circuitikz}}
\resposta{$I_1=\SI{3}{\ampere}$; $v_s=\SI{-15}{\volt}$}
\resolucao{Na malha 1,
$5I_1+5(I_1-2)=20$, então $I_1=3$ A. Na malha 2, percorrida no sentido horário,
\[
10(2)+v_s-5(3-2)=0,
\]
portanto $v_s=-15$ V. Isso significa que o terminal inferior da fonte está
$15$ V acima do terminal superior.}
\end{questao}

% =====================================================================
% N3 --- aprofundamento com novas topologias/objetivos
% =====================================================================
\blocoaprofundamento

\begin{questao}[3]{}
A fonte de corrente de $\SI{2}{\ampere}$ está entre as malhas 1 e 2. A fonte de
$\SI{8}{\volt}$ da malha 3 se opõe ao sentido horário indicado. Resolva usando
uma supermalha 1--2.
\circuitoq{%
\begin{circuitikz}[american,scale=.82,transform shape,line width=.8pt]
\draw (0,0) to[\fonteV,l^={$\SI{24}{\volt}$}] (0,3.4)
      to[R,l={$\SI{4}{\ohm}$}] (2.8,3.4) coordinate(t1);
\draw (t1) to[isource,l={$\SI{2}{\ampere}$},invert] (2.8,0) coordinate(b1)--(0,0);
\draw (t1) to[R,l={$\SI{2}{\ohm}$}] (5.6,3.4) coordinate(t2);
\draw (t2) to[R,l_={$\SI{4}{\ohm}$}] (5.6,0) coordinate(b2)--(b1);
\draw (t2) to[R,l={$\SI{4}{\ohm}$}] (8.4,3.4)
      to[\fonteV,l_={$\SI{8}{\volt}$}] (8.4,0)--(b2);
\fvMeshCW{1.25,1.05}{.34}{$I_1$}
\fvMeshCW{4.15,1.05}{.34}{$I_2$}
\fvMeshCW{7.05,1.05}{.34}{$I_3$}
\draw[dashed,laranja,line width=.9pt,rounded corners=7pt] (-.45,-.45) rectangle (6.05,3.85);
\node[laranja,fill=white,inner sep=1pt] at(2.8,-.72){\footnotesize supermalha 1--2};
\end{circuitikz}}
\resposta{$I_1=\SI{1.5}{\ampere}$; $I_2=\SI{3.5}{\ampere}$;
$I_3=\SI{0.75}{\ampere}$}
\resolucao{As equações são
\[
4I_1+6I_2-4I_3=24,
\qquad I_2-I_1=2,
\qquad -4I_2+8I_3=-8.
\]
Resolvendo, obtém-se $I_1=3/2$ A, $I_2=7/2$ A e $I_3=3/4$ A.}
\end{questao}

\begin{questao}[3]{}
O ramo comum é uma fonte ideal de $\SI{6}{\volt}$, positiva no topo. A fonte de
$\SI{18}{\volt}$ impulsiona $I_1$, enquanto a de $\SI{12}{\volt}$ se opõe a
$I_2$. Determine as correntes e a potência da fonte compartilhada.
\circuitoq{%
\begin{circuitikz}[american,scale=.92,transform shape,line width=.8pt]
\draw (0,0) to[\fonteV,l^={$\SI{18}{\volt}$}] (0,3.2)
      to[R,l={$\SI{6}{\ohm}$}] (3.2,3.2) coordinate(t);
\draw (t) to[\fonteV,l_={$\SI{6}{\volt}$}] (3.2,0) coordinate(b)--(0,0);
\draw (t) to[R,l={$\SI{4}{\ohm}$}] (6.4,3.2)
      to[\fonteV,l_={$\SI{12}{\volt}$}] (6.4,0)--(b);
\fvMeshCW{1.55,1.05}{.40}{$I_1$}
\fvMeshCW{4.85,1.05}{.40}{$I_2$}
\end{circuitikz}}
\resposta{$I_1=\SI{2}{\ampere}$; $I_2=\SI{-1.5}{\ampere}$;
$P_{6V}=\SI{21}{\watt}$ absorvidos}
\resolucao{Na malha 1, $6I_1+6=18$, logo $I_1=2$ A. Na malha 2,
$4I_2+12-6=0$, então $I_2=-1{,}5$ A. A corrente real na fonte comum, para
baixo, é $I_1-I_2=3{,}5$ A. Como ela entra no terminal positivo superior,
\[
P=6(3{,}5)=21\,\text{W},
\]
portanto a fonte de $6$ V absorve potência.}
\end{questao}

\begin{questao}[3]{}
As fontes de corrente periféricas fixam $I_1=\SI{2}{\ampere}$ e
$I_3=\SI{1}{\ampere}$. Determine $I_2$ e identifique qualquer ramo
compartilhado com corrente nula.
\circuitoq{%
\begin{circuitikz}[american,scale=.82,transform shape,line width=.8pt]
\draw (0,0) to[isource,l^={$\SI{2}{\ampere}$}] (0,3.4)
      --(2.8,3.4) coordinate(t1);
\draw (t1) to[R,l_={$\SI{2}{\ohm}$}] (2.8,0) coordinate(b1)--(0,0);
\draw (t1) to[R,l={$\SI{4}{\ohm}$}] (4.5,3.4)
      to[\fonteV,l={$\SI{11}{\volt}$}] (5.8,3.4) coordinate(t2);
\draw (t2) to[R,l_={$\SI{3}{\ohm}$}] (5.8,0) coordinate(b2)--(b1);
\draw (t2)--(8.6,3.4) to[isource,l={$\SI{1}{\ampere}$}] (8.6,0)--(b2);
\fvMeshCW{1.25,1.05}{.34}{$I_1$}
\fvMeshCW{4.30,1.05}{.34}{$I_2$}
\fvMeshCW{7.25,1.05}{.34}{$I_3$}
\end{circuitikz}}
\resposta{$I_2=\SI{2}{\ampere}$; o resistor de $\SI{2}{\ohm}$ conduz corrente nula}
\resolucao{A equação da malha central é
\[
(2+4+3)I_2-2I_1-3I_3=11.
\]
Com $I_1=2$ A e $I_3=1$ A:
$9I_2-4-3=11$, portanto $I_2=2$ A. Assim, no resistor compartilhado entre as
malhas 1 e 2, $I_1-I_2=0$.}
\end{questao}

\begin{questao}[3]{}
Analise a ponte resistiva pelo método das malhas e determine a corrente no
resistor de ponte de $\SI{5}{\ohm}$.
\circuitoq{%
\begin{circuitikz}[american,scale=.78,transform shape,line width=.8pt]
\coordinate (A) at (5,5); \coordinate (D) at (5,0); \coordinate (B) at (2.6,2.5); \coordinate (C) at (7.4,2.5);
\draw (A)--(0.5,5) to[\fonteV,l^={$\SI{18}{\volt}$}] (0.5,0)--(D);
\draw (A) to[R,l_={$\SI{2}{\ohm}$}] (B) to[R,l_={$\SI{4}{\ohm}$}] (D);
\draw (A) to[R,l={$\SI{3}{\ohm}$}] (C) to[R,l={$\SI{6}{\ohm}$}] (D);
\draw (B) to[R,l={$\SI{5}{\ohm}$}] (C);
\fvMeshCW{2.1,2.55}{.34}{$I_1$}
\fvMeshCW{5.0,3.55}{.32}{$I_2$}
\fvMeshCW{5.0,1.45}{.32}{$I_3$}
\end{circuitikz}}
\resposta{$I_1=\SI{5}{\ampere}$; $I_2=I_3=\SI{2}{\ampere}$;
$i_{ponte}=0$}
\resolucao{A matriz é
\[
\begin{bmatrix}
6&-2&-4\\
-2&10&-5\\
-4&-5&15
\end{bmatrix}
\begin{bmatrix}I_1\\I_2\\I_3\end{bmatrix}
=
\begin{bmatrix}18\\0\\0\end{bmatrix}.
\]
A solução é $(5,2,2)$ A. O resistor de ponte é comum às malhas 2 e 3, logo
$i_{ponte}=I_2-I_3=0$. O resultado expressa o equilíbrio da ponte.}
\end{questao}

\begin{questao}[3]{}
Use uma fonte de teste de $\SI{12}{\volt}$ e análise de malhas para determinar
a resistência equivalente vista pela fonte.
\circuitoq{%
\begin{circuitikz}[american,scale=.92,transform shape,line width=.8pt]
\draw (0,0) to[\fonteV,l^={$\SI{12}{\volt}$}] (0,3.2)
      to[R,l={$\SI{3}{\ohm}$}] (3.2,3.2) coordinate(t);
\draw (t) to[R,l_={$\SI{6}{\ohm}$}] (3.2,0) coordinate(b)--(0,0);
\draw (t) to[R,l={$\SI{6}{\ohm}$}] (6.4,3.2)--(6.4,0)--(b);
\fvMeshCW{1.55,1.05}{.40}{$I_1$}
\fvMeshCW{4.85,1.05}{.40}{$I_2$}
\end{circuitikz}}
\resposta{$I_1=\SI{2}{\ampere}$; $I_2=\SI{1}{\ampere}$;
$R_{eq}=\SI{6}{\ohm}$}
\resolucao{As equações são
\[
9I_1-6I_2=12,\qquad -6I_1+12I_2=0.
\]
Da segunda, $I_1=2I_2$; substituindo na primeira, $I_2=1$ A e $I_1=2$ A. A
fonte de teste fornece $2$ A, portanto $R_{eq}=12/2=6\,\ohm$.}
\end{questao}

% =====================================================================
% N4 --- sintese, otimizacao e problemas inversos
% =====================================================================
\blocodesafio

\begin{questao}[4]{}
A rede em ``roda'' possui quatro malhas triangulares em torno do nó $O$.
Monte a matriz e determine as quatro correntes.
\circuitoq{%
\begin{circuitikz}[american,scale=.72,transform shape,line width=.8pt]
\coordinate (A) at (-4,4); \coordinate (B) at (4,4); \coordinate (C) at (4,-4); \coordinate (D) at (-4,-4); \coordinate (O) at (0,0);
% lados externos: R=4 em serie com fontes
\draw (A) to[\fonteV,l^={$\SI{24}{\volt}$}] (-1.8,4) to[R,l={$\SI{4}{\ohm}$}] (B);
\draw (B) to[\fonteV,l={$\SI{12}{\volt}$}] (4,1.8) to[R,l={$\SI{4}{\ohm}$}] (C);
\draw (C) to[\fonteV,l_={$\SI{8}{\volt}$}] (1.8,-4) to[R,l_={$\SI{4}{\ohm}$}] (D);
\draw (D) to[\fonteV,l_={$\SI{4}{\volt}$},invert] (-4,1.8) to[R,l_={$\SI{4}{\ohm}$}] (A);
% raios
\draw (A) to[R,l={$\SI{2}{\ohm}$}] (O);
\draw (B) to[R,l_={$\SI{2}{\ohm}$}] (O);
\draw (C) to[R,l={$\SI{2}{\ohm}$}] (O);
\draw (D) to[R,l_={$\SI{2}{\ohm}$}] (O);
\fvMeshCW{0,2.45}{.34}{$I_1$}
\fvMeshCW{2.45,0}{.34}{$I_2$}
\fvMeshCW{0,-2.45}{.34}{$I_3$}
\fvMeshCW{-2.45,0}{.34}{$I_4$}
\node[fill=white,inner sep=1pt] at(O){$O$};
\end{circuitikz}}
\resposta{$I_1=\SI{4}{\ampere}$; $I_2=\SI{3}{\ampere}$;
$I_3=\SI{2}{\ampere}$; $I_4=\SI{1}{\ampere}$}
\resolucao{Cada malha contém $4\,\ohm$ exclusivo e dois raios de
$2\,\ohm$. A matriz é cíclica:
\[
\begin{bmatrix}
8&-2&0&-2\\
-2&8&-2&0\\
0&-2&8&-2\\
-2&0&-2&8
\end{bmatrix}
\begin{bmatrix}I_1\\I_2\\I_3\\I_4\end{bmatrix}
=
\begin{bmatrix}24\\12\\8\\-4\end{bmatrix}.
\]
A solução é $(4,3,2,1)$ A. O termo $-4$ representa a fonte da quarta malha
orientada contra a corrente horária de referência.}
\end{questao}

\begin{questao}[4]{}
Determine o resistor compartilhado $R$ para que $I_2=I_1/2$. Calcule também
as correntes resultantes.
\circuitoq{%
\begin{circuitikz}[american,scale=.92,transform shape,line width=.8pt]
\draw (0,0) to[\fonteV,l^={$\SI{20}{\volt}$}] (0,3.2)
      to[R,l={$\SI{4}{\ohm}$}] (3.2,3.2) coordinate(t);
\draw (t) to[R,l_={$R$}] (3.2,0) coordinate(b)--(0,0);
\draw (t) to[R,l={$\SI{6}{\ohm}$}] (6.4,3.2)--(6.4,0)--(b);
\fvMeshCW{1.55,1.05}{.40}{$I_1$}
\fvMeshCW{4.85,1.05}{.40}{$I_2$}
\end{circuitikz}}
\resposta{$R=\SI{6}{\ohm}$; $I_1=\dfrac{20}{7}\,\si{\ampere}$;
$I_2=\dfrac{10}{7}\,\si{\ampere}$}
\resolucao{Na malha 2,
\[
-RI_1+(6+R)I_2=0.
\]
Impondo $I_2=I_1/2$:
$-R+(6+R)/2=0$, logo $R=6\,\ohm$. Na malha 1,
$10I_1-6I_2=20$; com $I_2=I_1/2$, resulta $7I_1=20$.}
\end{questao}

\begin{questao}[4]{}
O resistor compartilhado é a carga $R_L$. Determine $R_L$ para que sua potência
seja máxima e calcule $P_{L,\max}$.
\circuitoq{%
\begin{circuitikz}[american,scale=.92,transform shape,line width=.8pt]
\draw (0,0) to[\fonteV,l^={$\SI{24}{\volt}$}] (0,3.2)
      to[R,l={$\SI{4}{\ohm}$}] (3.2,3.2) coordinate(t);
\draw (t) to[R,l_={$R_L$}] (3.2,0) coordinate(b)--(0,0);
\draw (t) to[R,l={$\SI{8}{\ohm}$}] (6.4,3.2)
      to[\fonteV,l_={$\SI{12}{\volt}$},invert] (6.4,0)--(b);
\fvMeshCW{1.55,1.05}{.40}{$I_1$}
\fvMeshCW{4.85,1.05}{.40}{$I_2$}
\end{circuitikz}}
\resposta{$R_L=\dfrac{8}{3}\,\ohm\approx\SI{2.67}{\ohm}$;
$P_{L,\max}=\SI{13.5}{\watt}$}
\resolucao{As equações são
\[
(4+R_L)I_1-R_LI_2=24,
\]
\[
-R_LI_1+(8+R_L)I_2=12.
\]
A corrente na carga é
\[
i_L=I_1-I_2=\frac{36}{3R_L+8}.
\]
Logo
\[
P_L=R_Li_L^2=\frac{1296R_L}{(3R_L+8)^2}.
\]
Derivando e igualando a zero, obtém-se $R_L=8/3\,\ohm$. Nesse ponto,
$i_L=36/16=2{,}25$ A e $P_{L,\max}=13{,}5$ W.}
\end{questao}

\begin{questao}[4]{}
A fonte $E$ deve ser escolhida para impor $I_2=0$. Determine $E$, sua polaridade
relativa ao sentido horário e as correntes $I_1$ e $I_3$.
\circuitoq{%
\begin{circuitikz}[american,scale=.82,transform shape,line width=.8pt]
\draw (0,0) to[\fonteV,l^={$\SI{24}{\volt}$}] (0,3.4)
      to[R,l={$\SI{6}{\ohm}$}] (2.8,3.4) coordinate(t1);
\draw (t1) to[R,l_={$\SI{2}{\ohm}$}] (2.8,0) coordinate(b1)--(0,0);
\draw (t1) to[R,l={$\SI{4}{\ohm}$}] (5.8,3.4) coordinate(t2);
\draw (t2) to[R,l_={$\SI{1}{\ohm}$}] (5.8,0) coordinate(b2)--(b1);
\draw (t2) to[R,l={$\SI{4}{\ohm}$}] (8.8,3.4)
      to[\fonteV,l_={$E$}] (8.8,0)--(b2);
\fvMeshCW{1.25,1.05}{.34}{$I_1$}
\fvMeshCW{4.30,1.05}{.34}{$I_2$}
\fvMeshCW{7.35,1.05}{.34}{$I_3$}
\end{circuitikz}}
\resposta{$I_1=\SI{3}{\ampere}$; $I_2=0$; $I_3=\SI{-6}{\ampere}$;
$E=\SI{30}{\volt}$ orientada contra $I_3$ horário}
\resolucao{O sistema é
\[
8I_1-2I_2=24,
\qquad
-2I_1+7I_2-I_3=0,
\qquad
-I_2+5I_3=E_s,
\]
onde $E_s$ é positivo se impulsiona $I_3$ horário. Impondo $I_2=0$:
$I_1=3$ A e, pela segunda equação, $I_3=-6$ A. Assim
$E_s=5(-6)=-30$ V. Portanto a fonte deve ter magnitude $30$ V e polaridade
oposta à referência horária.}
\end{questao}

\begin{questao}[4]{}
Em ensaio de uma rede de duas malhas, mediram-se $I_1=\SI{3}{\ampere}$ e
$I_2=\SI{2}{\ampere}$. Determine a resistência desconhecida $R$ do ramo comum
e a potência dissipada nela.
\circuitoq{%
\begin{circuitikz}[american,scale=.92,transform shape,line width=.8pt]
\draw (0,0) to[\fonteV,l^={$\SI{20}{\volt}$}] (0,3.2)
      to[R,l={$\SI{4}{\ohm}$}] (3.2,3.2) coordinate(t);
\draw (t) to[R,l_={$R$}] (3.2,0) coordinate(b)--(0,0);
\draw (t) to[R,l={$\SI{6}{\ohm}$}] (6.4,3.2)
      to[\fonteV,l_={$\SI{4}{\volt}$},invert] (6.4,0)--(b);
\fvMeshCW{1.55,1.05}{.40}{$I_1$}
\fvMeshCW{4.85,1.05}{.40}{$I_2$}
\end{circuitikz}}
\resposta{$R=\SI{8}{\ohm}$; $P_R=\SI{8}{\watt}$}
\resolucao{Pela malha 1,
\[
(4+R)3-R(2)=20 \Rightarrow 12+R=20 \Rightarrow R=8\,\ohm.
\]
A malha 2 confirma o mesmo valor:
\[
-3R+(6+R)2=4 \Rightarrow 12-R=4 \Rightarrow R=8\,\ohm.
\]
A corrente no ramo comum é $I_1-I_2=1$ A; portanto
$P_R=8(1)^2=8$ W.}
\end{questao}

% END BANCO COMPLEMENTAR

\gabarito[Capítulo 9 --- Método de Maxwell]
